\subsection{Analisis Revokasi dan Audit Delegasi}
\label{subsec:analisis-revokasi-audit-delegasi}

Delegasi akses menyebabkan token tidak lagi berdiri sendiri, tetapi membentuk relasi antara token induk dan token turunan. Oleh karena itu, mekanisme revokasi pada DecMed perlu diperluas agar pencabutan token induk dapat berdampak pada token turunan yang berasal dari token tersebut. Tanpa pencabutan berantai, aktor yang memperoleh token turunan masih dapat menggunakan akses meskipun akses asalnya telah dicabut oleh pasien.

Revokasi berantai diperlukan untuk mempertahankan kendali pasien sebagai pemilik data. Ketika pasien mencabut akses awal, sistem harus dapat mengenali bahwa token-token turunan yang berada dalam rantai delegasi yang sama tidak lagi memiliki dasar kewenangan yang sah. Dengan demikian, pemeriksaan revokasi tidak cukup hanya dilakukan terhadap token yang dikirim pada saat permintaan akses, tetapi juga perlu mempertimbangkan hubungan token tersebut dengan token asalnya.

Selain revokasi, delegasi juga membutuhkan audit yang mencatat hubungan antaraktor. Audit tidak cukup hanya mencatat bahwa suatu segmen diakses, tetapi juga perlu mencatat sumber kewenangan akses tersebut, termasuk aktor pemberi delegasi, aktor penerima delegasi, waktu delegasi, cakupan akses, dan token yang digunakan. Informasi ini diperlukan agar pasien dan sistem dapat menelusuri bagaimana akses terhadap data RME terbentuk selama proses pelayanan.

Berdasarkan analisis tersebut, pencabutan akses berantai dan audit log akses serta delegasi menjadi dasar bagi \texttt{S-TA-06}. Solusi ini melengkapi mekanisme delegasi \textit{offline} dengan kontrol setelah token diterbitkan, sehingga akses yang sudah tidak berlaku dapat ditolak dan riwayat delegasi tetap dapat ditelusuri.
